\section{Preliminaries}
\label{sec:foundations}

We collect the probabilistic and analytic identities used in the proof.
The Thorin representation reduces the problem to positive measures on the rate variable, while the gamma--Dirichlet and posterior identities will allow us to differentiate powers of the corresponding random variables.

For $\beta,b>0$, $\operatorname{Gamma}(\beta,b)$ denotes the gamma distribution with shape $\beta$ and rate $b$, having density
\[
 \frac{b^\beta}{\Gamma(\beta)}x^{\beta-1}e^{-bx},\qquad x>0,
 \qquad
 \Gamma(\beta)=\int_0^\infty x^{\beta-1}e^{-x}\,\dd x.
\]
We write $G_\beta$ for a random variable with distribution $\operatorname{Gamma}(\beta,1)$; thus $G_\beta/b$ has rate $b$.
The symbol $\perp$ denotes independence.

We use the integration notation introduced above also for unbounded integrable tests.
Expressions such as $F(y^2)$ and $F(|y|)$ stand for $\int y^2F(\dd y)$ and $\int |y|F(\dd y)$.
The point mass at $x$ is denoted by $\delta_x$, and $\mathbf1_A$ denotes the indicator of a set $A$.
Thus $T_*\nu(A)=\nu(T^{-1}(A))$ for a measurable map $T$.
We use $L_\nu(s)=\int e^{-sx}\nu(\dd x)$ for the Laplace transform of a probability on $[0,\infty)$; this is distinct from the law notation $\Law(X)$.

We write $\mathcal P(\R)$ for the Borel probability measures on $\R$ and $\Ptwo$ for those with finite second moment.
Weak, or narrow, convergence is denoted by $\Rightarrow$ and means convergence against bounded continuous functions.
Curves with values in $\Ptwo$ will be continuous for the narrow topology unless otherwise stated.
Unless specified otherwise, $L^p$ spaces on real intervals use Lebesgue measure.
We write $C_c^k(D)$ for the $k$-times continuously differentiable real functions with compact support in the open set $D$, and $\|f\|_\infty$ for the supremum norm on the indicated domain.
Finally, $r_+=r^+=\max\{r,0\}$, $r_-=r^-=\max\{-r,0\}$, and $\log_+x=\max\{\log x,0\}$ for $x>0$.

We also use the digamma function
\[
  \psi_0(r)=\frac{\Gamma'(r)}{\Gamma(r)},\qquad r>0,
\]

\subsection{Thorin measures}

A probability law \(\rho\) on \([0,\infty)\) belongs to \(\GGC\) if and only if its Laplace transform has the representation
\begin{equation}\label{eq:thorin-representation}
  L_\rho(s):=\int_{[0,\infty)}e^{-sx}\rho(\dd x)
  =\exp\left\{-as-\int_{(0,\infty)}
                  \log(1+s/b)\,U(\dd b)\right\},\qquad s\ge0,
\end{equation}
where \(a\ge0\) and \(U\) is a positive Borel measure satisfying
\begin{equation}\label{eq:thorin-admissibility}
  \int_{(0,\infty)}\log(1+1/b)\,U(\dd b)<\infty.
\end{equation}
The measure \(U\) is the rate-form Thorin measure and \(a\) is the drift.
We will first work with zero drift and finite Thorin mass; these are separate restrictions.
This is the Laplace-transform form of the representation in \cite[Section~3.1, p.~29, equations~(3.1.1)--(3.1.2)]{Bondesson1992}.
Condition \eqref{eq:thorin-admissibility} combines the usual integrability conditions at zero and infinity.
It also implies that \(U\) is finite on compact subintervals of \((0,\infty)\).

\begin{lemma}[Weak closure and finite-gamma approximation]
\label{lem:ggc-closure}
The class \(\GGC\) is closed under weak limits that are probability laws.
Every law in \(\GGC\) is the weak limit of laws of the form
\begin{equation}\label{eq:finite-gamma-approximation}
  \Law\left(\sum_{j=1}^{m_n}\frac{G_{\beta_{n,j}}}{b_{n,j}}\right),
  \qquad \beta_{n,j}>0,\quad b_{n,j}>0,
\end{equation}
where the gamma variables in each sum are independent.
Thus the approximating laws may be chosen with zero drift and finite atomic Thorin measures, even when the limiting law has positive drift or an infinite Thorin measure.
\end{lemma}

\begin{proof}
The weak closure assertion is \cite[Theorem~3.1.5, pp.~34--35]{Bondesson1992}, applied to a limit that is a probability measure.
The last paragraph on p.~35 of the same reference states the approximation assertion explicitly, with approximating drifts zero and approximating Thorin measures having finitely many atoms.
The transform of a finite atomic measure \(U_n=\sum_j\beta_{n,j}\delta_{b_{n,j}}\) is the transform of the sum in \eqref{eq:finite-gamma-approximation}.
If the target law is \(\delta_0\), a single gamma variable with a rate tending to infinity also supplies the asserted approximation.
\end{proof}

For a law with representation \eqref{eq:thorin-representation}, differentiation at \(s>0\) gives
\begin{equation}\label{eq:ggc-log-derivative}
  -\partial_s\log L_\rho(s)
  =a+\int_{(0,\infty)}\frac{U(\dd b)}{s+b}.
\end{equation}
The differentiation follows from \eqref{eq:thorin-admissibility}, locally uniformly in positive \(s\).

\subsection{Dirichlet processes and gamma means}

The following convention for Dirichlet processes is the one used in \cite[Section~1, reprint p.~2]{James2005}.

\begin{definition}[Dirichlet processes]
\label{def:dirichlet-process}
For a nonzero finite measure \(V\) on a Borel space \(E\), \(P\sim\DP(V)\) means that \(P\) is a \emph{Dirichlet process}: it is a random probability measure whose masses on any measurable finite partition of \(E\) have the Dirichlet distribution with parameters equal to the corresponding \(V\)-masses.
For positive parameters \((\beta_1,\ldots,\beta_m)\), this Dirichlet distribution is the law of \((G_{\beta_i}/\sum_jG_{\beta_j})_{i=1}^m\), with independent gamma variables.
A cell of zero base mass has zero random mass almost surely and is omitted from this vector.
We use this notation both for rates, with \(E=(0,\infty)\), and for log rates, with \(E=\R\).
The base measure may be atomic, nonatomic, or mixed.
\end{definition}

\begin{lemma}[Gamma--Dirichlet normalization and Markov--Krein identity]
\label{lem:gamma-dirichlet}
Let \(U\) be a nonzero finite positive measure on \((0,\infty)\), with total mass \(B\).
Let \(P\sim\DP(U)\), and let \(G_B\) be independent of \(P\).
Then the random measure \(G_BP\) is the gamma process with shape measure \(U\) and unit rate: its masses on disjoint sets are independent gamma variables with shapes equal to their \(U\)-masses, and sets of zero \(U\)-mass receive zero mass.
In particular, if \(r\ge0\) is Borel measurable and
\[
  \int\log(1+r(b))\,U(\dd b)<\infty,
\]
then \(P(r):=\int r(b)P(\dd b)<\infty\) almost surely, and
\begin{equation}\label{eq:markov-krein}
  \E e^{-zG_BP(r)}
  =\E(1+zP(r))^{-B}
  =\exp\left\{-\int\log(1+zr(b))\,U(\dd b)\right\},
  \qquad z\ge0.
\end{equation}
Consequently, if \(U\) satisfies \eqref{eq:thorin-admissibility}, the zero-drift GGC \(X\) with Thorin measure \(U\) admits the representation
\begin{equation}\label{eq:gamma-dirichlet-untilted}
  X\ \overset{\mathrm d}=\ G_B\int b^{-1}P(\dd b).
\end{equation}
For \(s>0\), let \(X_s\) denote the exponential tilt of \(X\):
\[
  \Law(X_s)(\dd x)
  =\frac{e^{-sx}}{\E e^{-sX}}\,\Law(X)(\dd x).
\]
Then, with
\[
  M_P(s):=\int_{(0,\infty)}\frac{P(\dd b)}{s+b},
\]
one has
\begin{equation}\label{eq:gamma-dirichlet-tilt}
  X_s\ \overset{\mathrm d}=\ G_BM_P(s),\qquad G_B\ \text{independent of }P.
\end{equation}
The base law of \(P\) on the right is \(\DP(U)\) for every \(s\).
\end{lemma}

\begin{proof}
The normalization, independence, logarithmic condition, and identity \eqref{eq:markov-krein} are recorded in \cite[Section~1, reprint p.~2, equations~(1)--(3)]{James2005}.
We recall the finite-partition argument to make explicit its extension to the unbounded functions used here.
For a finite partition with positive parameter vector \((\beta_1,\ldots,\beta_m)\), the change of variables \(x_j=tp_j\), with Jacobian \(t^{m-1}\), factors the joint density of independent unit-rate gamma variables into a \(\operatorname{Gamma}(B,1)\) density in \(t\) and a Dirichlet density in \((p_1,\ldots,p_m)\).
Conversely, multiplying an independent Dirichlet vector by \(G_B\) gives these independent gamma variables.
The gamma-process assertion follows by applying this calculation to every finite partition.

The Laplace functional for a nonnegative simple function is therefore the right side of \eqref{eq:markov-krein}.
Increasing simple approximation proves the same identity for an arbitrary nonnegative Borel function, initially allowing an infinite integral.
Under the stated logarithmic condition, the right side tends to one as \(z\downarrow0\), by dominated convergence with \(0\le z\le1\).
It follows that \(G_BP(r)\), and hence \(P(r)\), is finite almost surely.
Conditioning on \(P\) and integrating the gamma density gives the middle expression in \eqref{eq:markov-krein}.

Taking \(r(b)=b^{-1}\) proves \eqref{eq:gamma-dirichlet-untilted}.
For the tilted identity, \(r(b)=(s+b)^{-1}\) is bounded and
\begin{align*}
  \E e^{-vG_BM_P(s)}
  &=\exp\left\{-\int
       \log\left(1+\frac{v}{s+b}\right)U(\dd b)\right\}\\
  &=\frac{\E e^{-(s+v)X}}{\E e^{-sX}},\qquad v\ge0.
\end{align*}
Uniqueness of Laplace transforms proves \eqref{eq:gamma-dirichlet-tilt}.
The same auxiliary pair \((G_B,P)\) can thus represent each tilted marginal; only these marginal identities will be used.
\end{proof}

For later use, \(\operatorname{Beta}(1,B)\) denotes the distribution on \((0,1)\) with density \(B(1-z)^{B-1}\).
Sampling one point from a Dirichlet process adds a unit atom to its base measure in the posterior law.
The following form will turn random rate integrals into deterministic ones.

\begin{lemma}[One-observation posterior and Palm identity]
\label{lem:dirichlet-posterior}
Let \(U\) be a nonzero finite positive measure on \((0,\infty)\) with total mass \(B\).
For every nonnegative jointly Borel function \(\Phi(b,P)\),
\begin{equation}\label{eq:dirichlet-palm}
  \E_{\DP(U)}\int\Phi(b,P)P(\dd b)
  =\frac1B\int U(\dd b)\,
          \E_{\DP(U+\delta_b)}\Phi(b,P).
\end{equation}
The identity also holds for a signed jointly Borel function whenever \(\E_{\DP(U)}\int|\Phi(b,P)|P(\dd b)<\infty\).
A jointly measurable realization of the posterior laws is
\begin{equation}\label{eq:palm-coupling}
  P^{(b)}=(1-Z)Q+Z\delta_b,\qquad
  Q\sim\DP(U),\quad Z\sim\operatorname{Beta}(1,B),
  \qquad Q\ \text{independent of }Z,
\end{equation}
for which \(P^{(b)}\sim\DP(U+\delta_b)\).
\end{lemma}

\begin{proof}
The posterior \(P\mid Y=b\), when \(Y\mid P\) has law \(P\), is \(\DP(U+\delta_b)\); see \cite[Section~2, reprint pp.~4--5, including equation~(8)]{James2005}.
We give a proof using finite partitions.
Multiplication of a Dirichlet density with parameter vector \((\beta_j)_j\) by its \(i\)-th coordinate changes the parameter to \((\beta_j+\mathbf1_{\{j=i\}})_j\), with multiplier \(\beta_i/B\).
The assertion is trivial for a partition with one positive-mass cell.
On a common finite partition this proves \eqref{eq:dirichlet-palm} for cylinder functions and simple functions of the sampled location.
Refinement and the monotone-class theorem extend the identity to the stated jointly Borel functions.
Zero-mass cells contribute nothing.
This proof does not assume that \(U\) is diffuse.

To obtain \eqref{eq:palm-coupling}, add an independent \(\operatorname{Gamma}(1,1)\) mass at \(b\) to the gamma process \(G_BQ\), and normalize its total mass.
The new shape measure is \(U+\delta_b\).
The fraction assigned by the added mass is \(\operatorname{Beta}(1,B)\), independently of \(Q\).
The displayed formula, as a function of \((b,Q,Z)\), is a Borel map into the space of probability measures endowed with its weak Borel sigma-field.
It therefore also specifies the required measurable posterior kernel.
\end{proof}

Atomic base measures are allowed throughout, so the identity applies in particular to finite gamma convolutions.

\begin{lemma}[Stick-breaking realization]\label{lem:stick-breaking}
Let \(B>0\), let \(H\) be a probability on \(\R\), let \((Y_j)_{j\ge1}\) be independent with law \(H\), and let \((V_j)_{j\ge1}\) be independent \(\operatorname{Beta}(1,B)\) variables, independently of the locations.
Then
\begin{equation}\label{eq:stick-breaking}
  Q=\sum_{j\ge1}W_j\delta_{Y_j},\qquad
  W_j=V_j\prod_{i<j}(1-V_i),
\end{equation}
is a probability measure almost surely and has law \(\DP(BH)\).
\end{lemma}

\begin{proof}
This is Sethuraman's construction \cite[Section~2, pp.~642--643, equation~(2.1), and Theorem~3.4, p.~645]{Sethuraman1994}.
Note that the remaining mass after \(m\) terms is \(\prod_{j\le m}(1-V_j)\); it decreases to zero almost surely because its expectation is \((B/(B+1))^m\).

For a finite partition with probabilities \(p_i\), a Dirichlet vector \(D\) with parameters \(Bp_i\) satisfies
\[
  D\ \overset{\mathrm d}=\ V e_J+(1-V)D',
\]
where \(e_i\) is the \(i\)-th coordinate unit vector, \(D'\) is an independent copy of \(D\), \(V\) has the displayed beta law, and \(J\) is independent with probabilities \(p_i\).
Indeed, conditional on \(J=i\), gamma addition gives the Dirichlet parameters \(Bp_j+\mathbf1_{\{j=i\}}\).
Its density relative to that of \(D\) is \(x_i/p_i\); mixing with weights \(p_i\) restores the original density.
Again zero-probability cells are omitted and the one-cell case is immediate.
This fixed point is unique: two initial probability vectors driven by the same \(V_j,J_j\) have \(\ell^1\) distance at most \(2\prod_{j\le m}(1-V_j)\) after \(m\) iterations.
The partition vector defined by \eqref{eq:stick-breaking} satisfies this fixed-point equation and hence is Dirichlet.
This proves the lemma on all finite partitions, including for atomic \(H\).
\end{proof}

When the parameters vary, these random measures can be realized on a single probability space.
The precise construction is given in Lemma~\ref{lem:parameterized-dirichlet} of Appendix~\ref{sec:measurable-realizations}.
It will be used both for measurability of the coefficients and for their weak continuity.

\subsection{The bounded Stieltjes phase}

The phase representation below is the analytic input that makes the jump kernel nonnegative.
We use the convention for Stieltjes functions in \cite[Chapter~2]{SSV2010}.

\begin{definition}[Stieltjes functions]
\label{def:stieltjes}
A Stieltjes function is a function \(S:(0,\infty)\to[0,\infty)\) with a representation
\[
  S(s)=\frac{c}{s}+d+\int_{(0,\infty)}
                         \frac{\sigma(\dd t)}{s+t},
  \qquad c,d\ge0,\qquad
  \int\frac{\sigma(\dd t)}{1+t}<\infty,
\]
where \(\sigma\) is a positive Borel measure.
\end{definition}

The bounded phase representation used below is a consequence of two results in \citet[Theorems~6.10 and~7.3]{SSV2010}: if \(S\) is a nonzero Stieltjes function, its reciprocal \(f=1/S\) has the representation
\begin{equation}\label{eq:ssv-exponential}
  f(s)=\exp\left\{c_0+\int_0^\infty
     \left(\frac{t}{1+t^2}-\frac1{s+t}\right)
                           \xi(t)\,\dd t\right\},
  \qquad 0\le\xi\le1,
\end{equation}
with \(c_0\in\R\) and \(\xi\) unique up to Lebesgue-null sets \cite[Theorem~6.10, pp.~58--59, and Theorem~7.3, p.~63]{SSV2010}.
In the terminology of that reference, \(1/S\) is a complete Bernstein function.
We need only the displayed representation and its uniqueness.
In particular, they apply to the reciprocal of every resolvent mean considered below.

\begin{lemma}[Measurable phase and anchor-one representation]
\label{lem:phase}
For every probability measure \(P\) on \((0,\infty)\), extend \(M_P\) to the slit plane by
\[
  M_P(z)=\int\frac{P(\dd b)}{z+b},
  \qquad z\in\mathbb C\setminus(-\infty,0].
\]
There is a jointly Borel function \((P,t)\mapsto\xi_P(t)\in[0,1]\) such that
\begin{equation}\label{eq:phase-anchor}
  \log M_P(s)-\log M_P(1)
  =\int_0^\infty\xi_P(t)
       \left(\frac1{s+t}-\frac1{1+t}\right)\dd t,
  \qquad s>0.
\end{equation}
Here the probability measures \(P\) carry the Borel sigma-field of the narrow topology.
For each fixed \(P\), the density in this representation is unique Lebesgue-almost everywhere.
A specific jointly Borel choice is
\begin{equation}\label{eq:phase-representative}
  \xi_P(t)=\limsup_{n\to\infty}
       \frac{\arg(1/M_P(-t+i/n))}{\pi}
       =\limsup_{n\to\infty}
       \frac{-\arg M_P(-t+i/n)}{\pi},\qquad t>0,
\end{equation}
where \(n\) runs through the positive integers and \(\arg\) is the principal argument.
The ordinary limit in \eqref{eq:phase-representative} exists for almost every \(t\).
Both \eqref{eq:phase-anchor} and all its positive-order \(s\)-derivatives are absolutely convergent, locally uniformly in \(s>0\), uniformly over \(P\).
\end{lemma}

The representation follows by subtracting the logarithmic formula \eqref{eq:ssv-exponential} at $s$ and at $1$.
The measurable boundary version and the absolute convergence assertions are proved in Appendix~\ref{sec:phase-proof}.

Normalization at \(1\) is useful here because it requires neither \(M_P(0+)<\infty\) nor a logarithmic moment of the large rates.
The extra assumption for the zero-anchored variant in \cite[Remark~6.11, p.~60]{SSV2010} is therefore not used.
All references to \cite{SSV2010} use the 2010 first edition with the authors' posted corrections.
